\section{最大似然估计}

\subsection{基本框架}

最大似然估计（Maximum Likelihood Estimation, MLE）是训练显式密度生成模型的通用方法。

\begin{figure}[H]
\centering
\includegraphics[width=0.95\textwidth]{slides-images/slide-017.jpg}
\caption{最大似然估计：用神经网络参数化密度函数 $p_W(X)$，通过最大化训练数据的对数似然来学习参数 $W$\protect\footnotemark}
\end{figure}
\footnotetext{Slides 第18页。}

假设训练数据 $\{x^{(1)}, x^{(2)}, \ldots, x^{(N)}\}$ 独立同分布（i.i.d.）地来自某个未知的真实分布 $p_{\text{data}}$。我们用神经网络参数化一个密度函数 $p_W(X)$，目标是找到使训练数据最可能出现的参数 $W$：

$$W^* = \arg\max_W \prod_{i=1}^{N} p_W(x^{(i)})$$

通过取对数（log trick），将乘积转化为求和：

$$W^* = \arg\max_W \sum_{i=1}^{N} \log p_W(x^{(i)})$$

\begin{knowledgebox}{似然（Likelihood）vs 概率（Probability）}
\textbf{概率}：固定分布，改变数据点，观察各数据点的概率。\\
\textbf{似然}：固定数据点，改变分布（参数），观察哪个分布让这些数据点最可能出现。\\
最大似然估计本质上是在所有可能的分布中，选择让已观测数据最「合理」的那个分布。
\end{knowledgebox}

\subsection{本章小结}

最大似然估计通过最大化训练数据在模型分布下的概率来学习模型参数。log trick 将乘积转化为求和，便于优化。MLE 是自回归模型和 VAE 的共同理论基础。

%% ========== Part 4: 自回归模型 ==========

